\documentclass[11pt]{article}

\usepackage[margin=1in]{geometry}
\usepackage{amsmath}
\usepackage{booktabs}
\usepackage[dvipsnames]{xcolor}
\usepackage[colorlinks=true,linkcolor=NavyBlue,urlcolor=NavyBlue]{hyperref}
\usepackage{enumitem}
\usepackage{caption}
\usepackage{fancyhdr}
\usepackage{titlesec}

\captionsetup{font=small,labelfont=bf}
\setlist{itemsep=2pt,parsep=0pt,topsep=3pt}
\titlespacing*{\section}{0pt}{14pt}{6pt}
\titlespacing*{\subsection}{0pt}{10pt}{4pt}

\pagestyle{fancy}
\fancyhf{}
\fancyhead[L]{\small STAT 638 Project 17 --- frequentist counterpart}
\fancyhead[R]{\small\thepage}
\renewcommand{\headrulewidth}{0.4pt}

\newcommand{\kwh}{\,\mathrm{kWh}}
\newcommand{\code}[1]{\texttt{\small #1}}

\title{\vspace{-1.2cm}\textbf{A Frequentist Treatment of the Same Project}\\[2pt]
\large Household electricity demand, estimated by maximum likelihood}
\author{Project 17}
\date{\today}

\begin{document}
\maketitle
\thispagestyle{fancy}

\begin{abstract}
\noindent
The assigned analysis is Bayesian. This note keeps the same household, the same
daily series, the same regression, and the same six questions, and states how
each answer is obtained from a maximum-likelihood fit instead of a posterior.
The headline specification is unchanged: two annual harmonics, a linear trend,
day-of-week effects, and an AR(1) error with Student-$t$ innovations, on the
kWh scale. With $1{,}417$ days and a weak prior, the point estimates sit next
to the posterior means already reported. What changes is the language of the
claim: a confidence interval and a $p$-value, rather than a posterior
probability.
\end{abstract}

\section{What stays the same}

The data preparation does not move. One-minute power from the UCI
\emph{Individual Household Electric Power Consumption} series is converted to
daily energy by summing \code{Global\_active\_power} and dividing by $60$. A
day is kept only when at least 95\% of its $1{,}440$ minutes are observed.
That leaves $1{,}417$ usable days on a gap-free calendar of $T = 1{,}440$
days, from 17 December 2006 through 25 November 2010, with $23$ days retained
as missing so that the autoregression is defined on every date.

The mean structure is the one the brief requires. For a response $z_t = g(y_t)$,
with $g$ either the identity or the logarithm,
\begin{align}
z_t &= \mathbf{x}_t^{\top}\boldsymbol\beta + \varepsilon_t, \label{eq:mean}\\
\mathbf{x}_t^{\top}\boldsymbol\beta
 &= \beta_0
  + \beta_1 \frac{t - \bar t}{365}
  + \sum_{j=1}^{J}\left[a_j \sin\frac{2\pi j t}{365} + b_j \cos\frac{2\pi j t}{365}\right]
  + \sum_{k=2}^{7}\delta_k \mathbf{1}\{\mathrm{dow}_t = k\}, \label{eq:design}\\
\varepsilon_t &= \phi\,\varepsilon_{t-1} + \eta_t,
\qquad
\eta_t \sim t_\nu(0,\sigma^2). \label{eq:ar}
\end{align}
Monday is the baseline day. The trend is centred at $\bar t$, so $\beta_0$ is
the level at the midpoint of the record. The harmonics use raw $t$, as
mandated. The specification that the holdout selected, called M7 in the main
report, is $J = 2$, AR(1), $t$ innovations, and $g$ the identity, so $z_t$ is
daily consumption in kWh.

The twelve-month holdout is also unchanged. Models are fitted on data through
30 November 2009 ($1{,}072$ observed days) and scored on the final $345$
observed days. Log score, CRPS, and empirical coverage do not depend on a
posterior. They remain the right way to decide whether a specification
forecasts.

\section{What the fit is}

There is no prior. The parameters
$(\boldsymbol\beta, \phi, \sigma, \nu)$ are fixed unknown constants, and the
fit is the value that maximises the likelihood of the observed days.

For Gaussian innovations the conditional likelihood, after the AR(1)
quasi-difference
$\tilde z_t = z_t - \phi z_{t-1}$ and
$\tilde{\mathbf{x}}_t = \mathbf{x}_t - \phi \mathbf{x}_{t-1}$,
is ordinary weighted least squares in $\boldsymbol\beta$. In R the Gaussian
special case is
\begin{center}
\code{arima(y, order = c(1, 0, 0), xreg = X, method = "ML")}.
\end{center}
The exact likelihood adds the stationary density of the first residual,
with variance $\sigma^2/(1-\phi^2)$. That is the same initial-error term the
Gibbs sampler uses; here it enters the likelihood once, at the maximum, rather
than at every draw.

Student-$t$ innovations use the same normal scale mixture as the sampler,
\[
\eta_t \mid \lambda_t \sim N\!\left(0,\ \sigma^2/\lambda_t\right),
\qquad
\lambda_t \sim \mathrm{Gamma}(\nu/2,\ \nu/2),
\]
but as an EM algorithm.
\begin{itemize}
\item \textbf{E-step.} At the current parameters,
$\hat\lambda_t = (\nu+1)/(\nu + \eta_t^2/\sigma^2)$.
A small weight is a day the likelihood downweights. These weights are the
conditional expectations of the latent scales, evaluated once at the maximum
likelihood estimate.
\item \textbf{M-step.} Weighted least squares for $\boldsymbol\beta$, then
updates for $\sigma^2$, $\phi$, and $\nu$.
\end{itemize}
The algorithm stops at one point
$(\hat{\boldsymbol\beta},\hat\phi,\hat\sigma,\hat\nu)$ and one set of weights
$\hat\lambda_t$.

The $23$ missing days are integrated out of the likelihood, by a Kalman filter
or by EM imputation of their conditional expectations. A parametric bootstrap
then carries the uncertainty in those days into the standard errors. The
Bayesian fit imputes them inside the sampler, so that uncertainty is already
inside every posterior draw.

Standard errors come from the observed Fisher information at the maximum
likelihood estimate. A function of the coefficients --- an amplitude, a phase,
the weekend contrast, a percent-per-year slope, the horizon factor
$1/\sqrt{1-\phi^2}$ --- is handled by the delta method. When that approximation
is poor, as it is for an amplitude near zero or a phase near a branch cut, the
interval is read off a parametric bootstrap instead.

\section{The six questions}

The numerical anchors below are the posterior summaries already computed for
M7. Under a weak prior they are a close guide to the maximum likelihood
estimates. Each paragraph says what the frequentist report would quote.

\subsection{Q1 --- variation over the study period}

The midpoint level, the seasonal range, and the one-step innovation standard
deviation are maximum likelihood estimates, with confidence intervals from the
information matrix or the bootstrap. In the Bayesian fit those figures are a
midpoint level of $25.96\kwh$/day, a full seasonal range of $18.25\kwh$/day,
a weekly contribution of about $4.56\kwh$, and an innovation standard
deviation of $5.59\kwh$. The ordering is the same one the exploratory $R^2$
decomposition already gave: the first harmonic carries almost all of the
systematic variance ($R^2$ increment $0.36$), day of week adds $0.05$, and
the trend adds $0.01$.

The Newey--West intervals in the exploratory note are a descriptive frequentist
summary that ignores the AR($t$) structure. The model-based analysis replaces
them with intervals from the fitted autoregression.

\subsection{Q2 --- long-term trend}

The slope estimate is $-0.186\kwh$/year. The 95\% interval already computed,
$[-0.667,\ +0.292]$, implies a standard error of about $0.245$, a $z$-statistic
of about $-0.76$, and a two-sided $p$-value for $H_0:\beta_1 = 0$ of about
$0.45$. The interval covers zero. Four years do not establish a trend, and the
implied change over the whole record is under $3\%$ of the mean level.

The Bayesian report also quotes
$\Pr(\beta_1 < 0 \mid \mathrm{data}) = 0.775$. That number is the posterior
mass below zero. With a flat prior and a normal sampling distribution it is
$\Phi(0.76) \approx 0.78$, which is why the two analyses look so similar. The
frequentist summary stops at the estimate, the confidence interval, and the
test: the point estimate is negative, and a test of a zero slope does not
reject at conventional levels.

\subsection{Q3 --- annual pattern, and the second harmonic}

Amplitude and peak date are
\[
\hat A_j = \sqrt{\hat a_j^2 + \hat b_j^2},
\qquad
\widehat{\mathrm{phase}}_j = \mathrm{atan2}(\hat a_j,\hat b_j),
\]
with delta-method or bootstrap intervals. In the Bayesian fit the first
harmonic has amplitude $8.08\kwh$ and peaks on 14 January; the second has
amplitude $2.43\kwh$, and the amplitude ratio is $0.30$.

Harmonic 2 against harmonic 1 is a likelihood-ratio test on two extra
coefficients, referred to $\chi^2_2$, together with AIC. The comparison is
nested and regular. The substantive reason to keep the second harmonic is
unchanged: it deepens and delays the summer trough and sharpens the winter
peak. The Bayesian statement
$\Pr(A_2 > 0.15\,A_1 \mid \mathrm{data}) = 0.999$
becomes a confidence interval for the ratio that lies well above $0.15$.

\subsection{Q4 --- weekly pattern}

Each day-of-week deviation has a Wald interval from
$\widehat{\mathrm{Var}}(\hat{\boldsymbol\beta})$. The weekend-minus-weekday
contrast is the linear combination $\mathbf{c}^{\top}\hat{\boldsymbol\beta}$,
with standard error
$\sqrt{\mathbf{c}^{\top}\widehat{\mathrm{Var}}(\hat{\boldsymbol\beta})\,\mathbf{c}}$.
The Bayesian fit puts Saturday and Sunday about $3.5\kwh$ and $3.0\kwh$ above
the weekly mean, and the weekend-minus-weekday contrast at $4.56\kwh$/day
with interval $[3.71,\ 5.42]$. The posterior probability that the weekend
exceeds the weekday is $1$. The frequentist report is the same contrast, a
confidence interval lying entirely above zero, and a very small $p$-value for
a zero contrast.

\subsection{Q5 --- accuracy, and how uncertainty grows with the horizon}

The $h$-step point forecast plugs the maximum likelihood estimate into the
same recursion the sampler uses,
\[
\hat z_{T+h}
 = \mathbf{x}_{T+h}^{\top}\hat{\boldsymbol\beta}
 + \hat\phi\,\hat\varepsilon_{T+h-1},
\]
with future errors propagated from the last observed residual. For Gaussian
innovations a plug-in prediction interval uses the scale
\[
\hat\sigma\sqrt{\sum_{j=0}^{h-1}\hat\phi^{2j}}.
\]
For $t$ innovations the interval uses the quantiles of $t_{\hat\nu}$ on that
same scale.

With $\hat\phi = 0.375$ the width rises by about 8\% and then stops. The
ceiling is the stationary factor $1/\sqrt{1-\hat\phi^2} \approx 1.08$. In the
Bayesian fit the 95\% predictive interval widens from $28.9\kwh$ at one day
ahead to about $31\kwh$ by day seven, and the empirical coverage is $94.5\%$
at $h = 1$ and $95.2\%$ at $h = 7$. The frequentist calculation reproduces
that shape. A plug-in interval treats
$(\hat{\boldsymbol\beta},\hat\phi,\hat\sigma,\hat\nu)$ as known, so it is
slightly narrow. A parametric bootstrap of whole forecast paths puts the
estimation error back in. Knowing yesterday's consumption still buys only a
few days of sharper intervals. Beyond about a week the forecast is the
seasonal-plus-weekly mean.

\subsection{Q6 --- periods the mean does not explain}

One-step anomalies are tails of the fitted innovation distribution at the
maximum likelihood estimate: a day is flagged when its conditional residual is
extreme under $t_{\hat\nu}(0,\hat\sigma^2)$. A one-step residual cannot see a
level shift that the AR(1) has already absorbed, so a sustained anomaly is
still judged by a block statistic. The same 14-day mean used in the main
report is referred to the distribution of that mean under the fitted
stationary process, by simulation from the maximum likelihood estimate or by
a parametric bootstrap.

The episodes do not change. The August blocks are the house running near a
base load through the French summer holiday, and the winter blocks are high
days the seasonal mean does not reach. Both are behavioural or
meteorological, and neither is a function of the calendar terms in
\eqref{eq:design}.

\section{Choosing among the seven specifications}

Nested Gaussian comparisons use likelihood-ratio tests and AIC. AR(2) against
AR(1) adds one coefficient. $J = 2$ against $J = 1$ adds two. The logarithm
against the identity is not a nested comparison; both maximised likelihoods
have to be densities for daily kWh, so the log-scale fit contributes an extra
$-\log y_t$. Gaussian innovations against $t$ innovations are compared by AIC,
because $\nu \to \infty$ sits on the boundary of the parameter space.

The holdout designs carry over as written. In the Bayesian comparison, WAIC
preferred AR(2) and the twelve-month forecast rejected it, because the extra
lag was fitting the absence blocks one step ahead. AIC can make the same
mistake. The forecasting question still selects M7: two harmonics, AR(1), $t$
innovations, kWh scale. The log transform remains the worse forecasting scale.

\section{How the reported numbers translate}

\begin{center}
\small
\begin{tabular}{@{}p{0.34\textwidth}p{0.58\textwidth}@{}}
\toprule
\textbf{Quantity in the Bayesian report} & \textbf{Frequentist counterpart} \\
\midrule
Posterior mean and 95\% credible interval
  & Maximum likelihood estimate and 95\% confidence interval \\
Slope $-0.19\kwh$/year, interval covering 0
  & Same estimate, same interval, two-sided $p \approx 0.45$ \\
$\Pr(\mathrm{slope}<0 \mid \mathrm{data}) = 0.775$
  & No matching probability statement; the interval and the test are the summary \\
Amplitude, weekend contrast, $\phi$
  & The same point estimates, with delta-method or bootstrap intervals \\
Predictive interval at horizon $h$
  & Plug-in interval, or a bootstrap interval that includes estimation error \\
Posterior predictive $p$-value
  & Tail probability at the fitted model, or a bootstrap $p$-value \\
WAIC
  & AIC, plus the same holdout log score, CRPS, and coverage \\
\bottomrule
\end{tabular}
\end{center}

\section{What the conclusions still say}

The annual cycle dominates. The second harmonic is needed because the summer
trough is deeper and later than a single sine wave. The weekend premium is
several kWh per day, and its confidence interval lies entirely above zero. A
linear trend is not resolved in four years. Forecast uncertainty stops
widening after about four days, because $\phi$ is only about $0.4$. Beyond a
week, a forecast for this house is seasonal-plus-weekly climatology.

The assignment asks for posterior probabilities, and those statements stay in
the main report. This note records the sampling-theory version of the same
fit, for the same questions, on the same series.

\end{document}
